% Nouveau EGA V (cote 25): archive cote 25, site page 175 (PDF page 176).
% typescript. Draft compared: none.
% AI-transcribed by Claude Opus 5.5 (claude-opus-5-5, adaptive thinking, effort high), Message Batches API, batch ega25b (msgbatch_01A7n5YuXubUPNK1zsepi8RG), 2026-11-13; not yet checked by eye.
\scanpage{25/175}
\mspage{\hand{2.}}

associée à Q (définie par la matrice des $\frac{\hand{\partial} Q}{\hand{\partial} X_i}$ ) est alternée, donc son
``noyau''  $N\subset\overset{\hand{\vee}}{E}$  est tel que  $\overset{\hand{\vee}}{E}/N$  soit de rang \uline{pair} , donc si E est lui-%
même de rang pair, il en est de même de N. Donc si $N\neq 0$, le rang de N
est au moins 2, d'où s'ensuit qu'il existe au moins un zéro non banal de
Q dans N, i.e. Q est non lisse.  NB Ce recours aux coordonnées est décidé
ment indigne de nous, il faudrait les utiliser uniquement pour prouver
simplement le
\uline{Lemme} \inserted[typed]{16.3.} E vectoriel de rang fini sur un corps k, \inserted[typed]{$Q\in \mathrm{Sym}^2(E)$} x un élément non nul de
$\overset{\hand{\vee}}{E}$ , \xout{\unsure{tel que Q(x}} x' son image dans P(E), supposons  $Q(x,x) = 0$ i.e.
$x'\in V(Q)$. Alors pour que \xout{\unsure{x' soit un point lisse de}} V(Q) soit lisse en x',
et de dim \xout{\illegible{1}} rang E - 1   (i.e. $Q\neq 0$), il f et s  que x n'appar-%
tienne pas au noyau de l'homomorphisme $\overset{\hand{\vee}}{E}\longrightarrow E$ défini par la forme bilinéai-%
re symétrique B associée à Q/
16.4.  L'étude qui suit est destinée essentiellement à \xout{\unsure{traiter le cas}}
donner un critère de lissité pour une forme quadratique dans le cas non
couvert par 16.2. , i.e. essentiellement dans le cas d'un vectoriel de
dimension impaire sur un corps de car 2.\xout{\unsure{Alors}} \inserted[typed]{Dans ce} cas \uline{toute} forme quadra-%
tique  Q  sur $\overset{\hand{\vee}}{E}$ est dégénérée, mais on voit aisément (prenant notamment
% NOTE: a line of typed text overtyped with x stands just below "est dégénérée," between this line and the next
\xout{\illegible{2}}
la``forme standart'' ) qu'elle peut fort bien être lisse.

Introduisons, pour tout entier  $n>0$, la ``forme quadratique standart''
$Q_n$ sur $\underline{Z}^n$ , comme une forme à coéfficients entiers en n variables
$X_1,..,X_n$ ,  en distinguant les deux cas:
a)  $n = 2m$  ,  $Q_{2m}(X_1,..,X_{2m}) = X_1X_2 + X_3X_4 + \ldots$
b)  $n = 2m+1$,  $Q_{2m+1}(X_1,..,X_{2m+1}) =$\xout{\unsure{$(X_1X_2+X$}} $Q_{2m}(X_1,..,X_{2m}) + X_{2m+1}^2$.
\uline{Lemme} 16.5  Soit n un entier impair, et considérons
\[ \xout{\unsure{M}}Q(X_{\xout{\unsure{2}}1},..,X_n) = \sum_i a_iX_i^{\,2} + \hand{\sum_{i<j}} b_{ij}X_iX_j \ , \]
une forme quadratique en n \unsure{indétermihées|indéterminées} $X_i$ , à coéff indéterminés
$a_i$, $b_{ij}$, de sorte que le discriminant  d(Q)  est un polynôme à

\mspage{\hand{175}}
